\documentclass[11pt,a4paper]{article}
\usepackage[margin=0.75in]{geometry}
\usepackage{enumitem}
\usepackage[hidelinks]{hyperref}
\setlist[itemize]{leftmargin=*,topsep=2pt,itemsep=1pt}
\pagestyle{empty}
\begin{document}

\begin{center}{\LARGE\bfseries AI / LLM Agent Engineer – Agentic LLM Systems with Evaluation and Guardrails}\\[4pt]
Aniket Saxena | h20250169@pilani.bits-pilani.ac.in | github.com/aniketqw\end{center}

\section*{Summary}
Built end‑to‑end data ingestion pipelines that clean, chunk, extract, validate, and store information in MongoDB, ensuring idempotent updates and source‑verified facts. Developed a LangChain ReAct Agent powered by Ollama, integrated with MongoDB for knowledge retrieval and structured output. Designed modular demo scripts to isolate and debug each pipeline stage, facilitating rapid iteration and testing.

\section*{Projects}
\begin{itemize}
  \item Built a data ingestion pipeline that processes data through acquisition, cleaning, chunking, extraction, validation, and storage in MongoDB.
  \item Implemented idempotent pipeline logic that updates existing records instead of duplicating them.
  \item Validated facts by ensuring the source quote appears verbatim in the originating chunk.
  \item Used an LLM exclusively for semantic extraction while keeping all other steps deterministic Python.
  \item Created demo scripts that isolate each pipeline stage for debugging.
  \item Implemented a LangChain ReAct Agent with Ollama and MongoDB for agentic LLM interactions.
\end{itemize}

\section*{Skills}
\begin{itemize}
  \item Python
  \item MongoDB
  \item LangChain
  \item Ollama
  \item LLM
  \item Semantic Extraction
  \item Data Validation
  \item Docker Compose
  \item Pipeline Debugging
\end{itemize}

\end{document}